\subsection{未来工作}
\label{sec:discussion-future}

从本研究可以得出几个具体方向。首先，训练协议依赖于 Ti-6Al-4V 试样的配对数据；将 PMDF-Net 扩展至碳纤维增强聚合物和复合层压板将需要域适应技术，例如对抗性微调或物理一致迁移学习，以考虑材料相关的 Lamb 波频散特性。其次，在 safety-critical 检测工作流程中部署需要置信度估计以及点预测；贝叶斯神经网络层或蒙特卡罗 dropout 可在推理时产生方差估计，标记分布外输入，使网络能够在不确定性超过操作阈值时弃权或触发人工审核。

第三，在 Jetson Orin 等便携式平台上进行实时现场部署需要量化网络权重和分析驱动的算子融合以满足延迟约束；在嵌入式硬件上的系统性基准测试将确定在代表性信号批处理下可达到的推理吞吐量。第四，当前架构仅通过损失函数嵌入物理约束；将 dispersive Lamb 波频散关系作为结构化操作硬连接到网络拓扑中可提高对新几何形状的泛化能力，而无需进行特定几何形状的 retraining。第五，配对数据需求限制了高平均参考信号可用的领域的适用性；基于 CycleGAN 或对抗性去噪的方法可实现从非配对噪声/干净语料库进行训练， substantially expanding 可到达的应用空间。最后，缺陷数据集主要覆盖切口和平底孔；扩展以包括疲劳裂纹、应力腐蚀开裂和分层将建立 PMDF-Net 在工业检测中遇到的各类结构相关缺陷类型上的性能。